\clearpage\fancyhead[L]{\small INSTRUMENTE.}\hypertarget{native-370-lower}{}\label{native-370-lower}
\begin{center}\textbf{INSTRUMENTE.}\par\vspace{4pt}\hrule\vspace{4pt}\fontsize{10}{12}\selectfont Göttingische gelehrte Anzeigen. Stück 75. Seite 745\;..\;752. 10. Mai 1813.\par\vspace{4pt}\hrule\vspace{4pt}\end{center}

\par\indent In der Sitzung der Königlichen Sozietät der Wissenschaften am 3. April, in der Professor \textsc{Eichhorn} den Vortrag hielt, berichtete Professor \textsc{Gauss} über zwei neue astronomische Instrumente. Die hiesige Sternwarte erhielt sie vor einigen Monaten aus der Werkstatt des berühmten Instrumentenbauers, Herrn Salinenrat \textsc{Reichenbach} in München. Professor \textsc{Gauss} legte zugleich die ersten damit angestellten Beobachtungen vor.

\par\indent Das \emph{erste} ist ein Repetitionskreis mit zwei Fernrohren und einem Durchmesser von 12 Zoll. An diesem Kunstwerk bewundert man die Feinheit und Genauigkeit der Teilung, die fast unglaubliche Empfindlichkeit der Libellen und die Vollkommenheit der Fernrohre ebenso wie die Präzision und Schönheit der Ausarbeitung aller einzelnen Teile des Instruments, die Leichtigkeit seiner vielen Bewegungen und die verschiedenen sinnreich angeordneten Vorrichtungen, mit denen der geniale Instrumentenbauer dem Beobachter die Arbeit erleichtert hat.

\par\indent Der Hauptkreis ist auf eingelegtem Silber unmittelbar von 5 zu 5 Minuten geteilt. Jeder der vier Nonien unterteilt einen solchen Abschnitt wiederum in 75 Teile, also von 4 zu 4 Sekunden. Die Nonien stehen im rechten Winkel zueinander und befinden sich in derselben Ebene wie der Gradbogen auf einem vollständigen zweiten Kreis. Dieser bewegt sich innerhalb des ersten so vollkommen konzentrisch, dass er ihn, obwohl mit bloßem Auge gar kein Zwischenraum wahrnehmbar ist, \SourceLeaf{371}nirgends berührt und sich daher mit größter Leichtigkeit bewegen lässt. Auch durch die absichtlich etwas schräg gestellten Mikroskope erkennt man den Zwischenraum nicht. Die Striche der Nonien scheinen unmittelbar an die Striche des Gradbogens anzustoßen. Das erhöht die Ablesegenauigkeit beträchtlich, ebenso wie die zweckmäßige Beleuchtung und die außerordentliche Feinheit und Gleichmäßigkeit der Teilstriche. Man kann ohne Weiteres 2, nötigenfalls sogar einzelne Sekunden schätzen. Die vier Nonien weichen nie um mehr als ein paar Sekunden voneinander ab. Das beweist sowohl das Fehlen jeder merklichen Exzentrizität als auch die unübertreffliche Genauigkeit der Teilung selbst.

\begingroup\rightskip=0pt plus 1em \par\indent Die drei zum Kreis gehörigen Libellen sind außerordentlich empfindlich. Die Hauptlibelle am hinteren Fernrohr gibt bei einer Neigung um eine Sekunde einen Ausschlag von mehr als einer Pariser Linie. Die äußerst geringe und doch gleichmäßige Krümmung im Innern der Glasröhre setzt einen Radius von mehr als 1400 Fuß voraus und ließ sich nur durch sehr kunstvolle und sorgfältige Bearbeitung erreichen. Um diese Krümmung zu erhalten, wurden die Glasröhren an beiden Enden mit genau eingeschliffenen Glasstöpseln, über die noch eine Blasenhaut gezogen ist, vollkommen verschlossen, ohne die Enden zuzuschmelzen.\par\endgroup

\par\indent Die beiden Fernrohre sind nur 16 Zoll lang und haben eine Öffnung von 15 Linien. Dennoch sind sie von ganz ausgezeichneter Güte, vollkommen achromatisch und vertragen dank ihrer außerordentlichen Präzision starke Vergrößerungen. Daher lässt sich mit großer Schärfe auf die feinen, im Brennpunkt eingezogenen Spinnfäden zielen. Das vordere Fernrohr hat ein prismatisches Okular, sodass das Auge bei der Höhenmessung immer durch horizontale Strahlen blickt. Hoch stehende Sterne lassen sich daher ebenso bequem beobachten wie tief stehende. Der Lichtverlust ist dabei fast unmerklich.

\par\indent Alle Schrauben für die Feineinstellung sind mit größter Präzision gearbeitet. Sie sprechen auf die leiseste Berührung sofort richtig an, und besondere Bremsschrauben verhindern auch künftig jedes Spiel. Ihre Gewinde sind so fein, dass auf einen Zoll mehr als 100 Gänge kommen.

\par\indent Da hier nicht der Ort für eine vollständige Beschreibung des Instruments ist, übergehen wir mehrere ebenso sinnreiche wie zweckmäßige Vorrichtungen, die den \textsc{Reichenbach}schen Kreisen eigen sind. Dazu gehören die Mittel, die Kreisebene ganz genau senkrecht zu stellen, Gegenstände leicht aufzufinden und ins Gesichtsfeld zu bringen und dergleichen.

\par\indent Mehrere Umstände verzögerten den Beginn der astronomischen Beobachtungen mit diesem Kreis bis Mitte März. Professor \textsc{Gauss} hat seither mit Unterstützung von Professor \textsc{Harding} beobachtet, der die Libelle einstellte. Als Beispiel teilen wir hier die Ergebnisse für die Polhöhe mit, die sich aus den beobachteten unteren Kulminationen des Polarsterns ergeben haben. Ihre gute Übereinstimmung bezeugt am besten die Vorzüglichkeit des Instruments.

\Needspace{13\baselineskip}\begin{center}Polhöhe der Göttinger Sternwarte aus Beobachtungen des Polarsterns bei der unteren Kulmination\end{center}
\begin{center}\fontsize{10}{12}\selectfont\setlength{\tabcolsep}{7pt}\renewcommand{\arraystretch}{1.15}\begin{tabular}{ll|c|r}
& 1813 & \shortstack{Anzahl der\\Beobachtungen} &\\\hline
März & 20 & 10 & $51^\circ\;31'\;54''86$\\
& 22 & 18 & $55.73$\\
& 26 & 18 & $57.41$\\
& 31 & 18 & $56.25$\\
April & 3 & 18 & $56.55$\\
& 7 & 32 & $55.20$\\
& 8 & 22 & $57.33$\\
\end{tabular}\end{center}

\par\indent\SourceLeaf{372}Das Mittel aus allen 136 Beobachtungen beträgt $51^\circ\,31'\,56''20$. Das äußerste Ergebnis weicht davon nur um $1''34$ ab. Für die Reduktion auf den Mittelpunkt der Sternwarte sind davon noch $0''16$ abzuziehen. Damit beträgt ihre Polhöhe $51^\circ\,31'\,56''04$, nur $2''$ mehr als von \textsc{Tob. Meyer} bestimmt. Dieses Ergebnis hängt noch von der Deklination des Polarsterns ab, für die die Bestimmung von Herrn von \textsc{Zach} zugrunde gelegt wurde. Sobald sich der Polarstern bei der oberen Kulmination beobachten lässt, wird es sich daher vielleicht noch ändern, aber gewiss nur sehr wenig. Hätte man für die Deklination das Mittel aus den Bestimmungen der Herren von \textsc{Zach}, \textsc{Oriani}, \textsc{Bouvard} und \textsc{Pond} zugrunde gelegt (deren Extremwerte nur um $\frac{3}{4}$ Sekunden voneinander abweichen), oder die Bestimmung von \textsc{Pond} allein, die mit diesem Mittel genau übereinstimmt, so wäre die Polhöhe um weitere $0''44$ kleiner ausgefallen.

\par\indent Das \emph{zweite} Instrument, in seiner Art ein ebenso bewundernswertes Meisterwerk, ist ein Repetitionstheodolit von 8 Zoll Durchmesser. Der Horizontalkreis ist unmittelbar von 10 zu 10 Minuten geteilt. Jeder der vier Nonien gibt $10''$ an, kleinere Teile lassen sich noch schätzen. Auch hier liegen die Nonien auf einer vollständigen Kreisscheibe innerhalb des geteilten Kreises und in derselben Ebene. Der innere Kreis oder Noniuskreis trägt zwei Stützen, auf denen das Hauptfernrohr von 12 Zoll Länge und 13 Linien Öffnung wie bei einem Passageninstrument an einer horizontalen Achse aufgehängt ist, die in zwei vollkommen gleiche zylindrische Stahlzapfen ausläuft. Wie bei einem Meridianfernrohr lässt sich seine Visierlinie durch Umlegen ganz genau senkrecht zu dieser Achse ausrichten. Eine an die Achse anzuhängende Libelle dient dazu, die Kreisebene waagerecht und die Achse parallel zu ihr zu stellen. Bei dieser Libelle ergibt 1 Sekunde einen Ausschlag von einer halben Linie. Ehe das Fernrohr den Horizontalkreis berührt, lässt es sich bis zu $40^\circ$ über und unter den Horizont neigen. Diese Neigung wird an einem eigenen Vertikalkreis von 5 Zoll Durchmesser gemessen, dessen Nonius einzelne Minuten angibt und an dem sich noch halbe oder Drittelminuten schätzen lassen. Das zweite, untere Fernrohr ist dem oberen ganz gleich, dient aber nur als Kontrollfernrohr, um den unveränderten Stand des Instruments zu prüfen. Es ist erstaunlich, wie genau sich mit einem so kleinen Instrument Winkel messen. lassen. Einzelmessungen geben die Winkel stets bis auf wenige Sekunden genau, und durch Repetition lassen sich ihre wahren Werte bis auf $1''$, höchstens $2''$, sicher bestimmen. Für den astronomischen Gebrauch ist besonders wichtig, wie leicht und genau man durch Beobachtung der Sonne oder der Fixsterne (von denen man diejenigen erster Größe leicht am hellen Tag sieht) die Azimute irdischer Gegenstände bestimmen kann. Sind diese für einen bestimmten Standort einmal genau ermittelt, können sie ihrerseits zur bequemen, von der Refraktion unabhängigen Zeitbestimmung dienen. Dieser Vorteil ist besonders in den Wintermonaten an einer Sternwarte ohne Passageninstrument hoch einzuschätzen.

\par\Needspace{10\baselineskip}\indent Besonders erfreulich ist, dass die schönen Fernrohre an diesen Instrumenten ganz deutschen Ursprungs sind. In ihrer Vollkommenheit stehen sie den besten englischen Fernrohren gleicher Abmessungen in nichts nach und übertreffen sie womöglich sogar. Das Flintglas für die Objektive wird in Benedictbeuern hergestellt. --- [Diesem Bericht hat \textsc{Gauss} eine Preisliste der im optischen Institut von \textsc{Reichenbach}, \textsc{Frauenhofer} und von \textsc{Utzschneider} gefertigten optischen Instrumente beigefügt, die in diese Ausgabe seiner Werke nicht aufgenommen wurde.]
